\documentclass{beamer}

\usepackage{amsmath,amsfonts,amssymb,bm}
\usepackage{quantikz}
\usepackage{graphicx}

\usetheme{Madrid}

\title{Adaptive Variational Quantum Eigensolver (ADAPT-VQE)}
\subtitle{Variational Quantum Algorithms and Many-Body Physics}
\author{}
\date{}

\begin{document}

\frame{\titlepage}

%================================================
\section{Motivation}
%================================================

\begin{frame}{Why Adaptive Ansätze?}
\begin{itemize}
\item Standard VQE relies on fixed ansatz
\item Problems:
\begin{itemize}
\item Overparameterization
\item Barren plateaus
\item Poor physical interpretability
\end{itemize}
\item Goal: construct ansatz \emph{adaptively}
\end{itemize}
\end{frame}

%------------------------------------------------

\begin{frame}{Physics Motivation}
\begin{itemize}
\item Ground state problem:
\[
H|\psi_0\rangle = E_0 |\psi_0\rangle
\]
\item Variational principle:
\[
E[\psi] = \frac{\langle \psi|H|\psi\rangle}{\langle \psi|\psi\rangle}
\]
\item Seek optimal parameterized wavefunction
\end{itemize}
\end{frame}

%================================================
\section{Standard VQE}
%================================================

\begin{frame}{Variational Quantum Eigensolver}
\[
|\psi(\theta)\rangle = U(\theta)|0\rangle
\]
\[
E(\theta) = \langle \psi(\theta)|H|\psi(\theta)\rangle
\]
\begin{itemize}
\item Classical optimizer minimizes $E(\theta)$
\end{itemize}
\end{frame}

%------------------------------------------------

\begin{frame}{Ansatz Challenge}
\begin{itemize}
\item Hardware-efficient ansatz:
\begin{itemize}
\item Shallow but unstructured
\end{itemize}
\item Problem-inspired ansatz (e.g. UCC):
\begin{itemize}
\item Physically motivated but deep
\end{itemize}
\item Trade-off motivates adaptive methods
\end{itemize}
\end{frame}

%================================================
\section{Coupled Cluster Connection}
%================================================

\begin{frame}{Coupled Cluster Ansatz}
\[
|\psi\rangle = e^{T} |\phi_0\rangle
\]
\[
T = T_1 + T_2 + \cdots
\]
\begin{itemize}
\item $T_k$ generates $k$-particle excitations
\end{itemize}
\end{frame}

%------------------------------------------------

\begin{frame}{Unitary Coupled Cluster}
\[
|\psi\rangle = e^{T - T^\dagger} |\phi_0\rangle
\]
\begin{itemize}
\item Hermitian generator
\item Suitable for quantum circuits
\end{itemize}
\end{frame}

%================================================
\section{ADAPT-VQE Idea}
%================================================

\begin{frame}{Core Idea}
\begin{itemize}
\item Build ansatz iteratively
\item Select operators based on gradient
\item Add most important operator at each step
\end{itemize}
\end{frame}

%------------------------------------------------

\begin{frame}{Operator Pool}
\begin{itemize}
\item Define pool:
\[
\{A_k\}
\]
\item Example:
\begin{itemize}
\item Fermionic excitations
\item Pauli strings
\end{itemize}
\end{itemize}
\end{frame}

%================================================
\section{Gradient Selection}
%================================================

\begin{frame}{Gradient Criterion}
\[
g_k = \left| \frac{\partial E}{\partial \theta_k} \right|
\]
\[
= \left| \langle \psi | [H, A_k] | \psi \rangle \right|
\]
\end{frame}

%------------------------------------------------

\begin{frame}{Physical Interpretation}
\begin{itemize}
\item Measures how strongly operator $A_k$ lowers energy
\item Equivalent to residual in many-body methods
\end{itemize}
\end{frame}

%================================================
\section{ADAPT-VQE Algorithm}
%================================================

\begin{frame}{Algorithm Steps}
\begin{enumerate}
\item Initialize reference state
\item Compute gradients for all $A_k$
\item Select operator with largest gradient
\item Add to ansatz
\item Re-optimize parameters
\item Repeat until convergence
\end{enumerate}
\end{frame}

%------------------------------------------------

\begin{frame}{Ansatz Growth}
\[
U(\theta) =
e^{\theta_n A_n} \cdots e^{\theta_1 A_1}
\]
\end{frame}

%================================================
\section{Circuit Representation}
%================================================

\begin{frame}{Single Operator Circuit}
\[
\begin{quantikz}
\lstick{|0\rangle} & \gate{e^{\theta A_k}} & \qw
\end{quantikz}
\]
\end{frame}

%------------------------------------------------

\begin{frame}{Layered Circuit}
\[
\begin{quantikz}
\lstick{|0\rangle} & \gate{A_1} & \gate{A_2} & \gate{A_3} & \qw
\end{quantikz}
\]
\end{frame}

%================================================
\section{Measurement of Gradients}
%================================================

\begin{frame}{Commutator Evaluation}
\[
\langle [H,A_k] \rangle
\]
\begin{itemize}
\item Requires decomposition into Pauli terms
\end{itemize}
\end{frame}

%------------------------------------------------

\begin{frame}{Measurement Cost}
\begin{itemize}
\item Scales with operator pool size
\item Trade-off:
\begin{itemize}
\item Smaller pool → faster
\item Larger pool → more expressive
\end{itemize}
\end{itemize}
\end{frame}

%================================================
\section{Convergence}
%================================================

\begin{frame}{Stopping Criterion}
\[
\max_k g_k < \epsilon
\]
\end{frame}

%------------------------------------------------

\begin{frame}{Energy Convergence}
\begin{itemize}
\item Typically exponential convergence
\item Often fewer parameters than UCC
\end{itemize}
\end{frame}

%================================================
\section{Advantages}
%================================================

\begin{frame}{Why ADAPT-VQE?}
\begin{itemize}
\item Compact ansatz
\item Problem-specific structure
\item Reduced circuit depth
\end{itemize}
\end{frame}

%------------------------------------------------

\begin{frame}{Physics Insight}
\begin{itemize}
\item Builds correlations step by step
\item Analogous to:
\begin{itemize}
\item Coupled cluster iterations
\item Many-body perturbation expansion
\end{itemize}
\end{itemize}
\end{frame}

%================================================
\section{Challenges}
%================================================

\begin{frame}{Main Limitations}
\begin{itemize}
\item Large measurement overhead
\item Operator pool design
\item Noise sensitivity
\end{itemize}
\end{frame}

%================================================
\section{Extensions}
%================================================

\begin{frame}{Variants}
\begin{itemize}
\item Qubit-ADAPT
\item Batched ADAPT
\item Gradient-free variants
\end{itemize}
\end{frame}

%================================================
\section{Summary}
%================================================

\begin{frame}{Summary}
\begin{itemize}
\item ADAPT-VQE constructs ansatz dynamically
\item Uses gradient to guide operator selection
\item Strong connection to many-body methods
\item Promising for NISQ devices
\end{itemize}
\end{frame}

%================================================
\section{Variational Principle and Many-Body Context}
%================================================

\begin{frame}{Ground State Problem}
\[
H|\psi_0\rangle = E_0 |\psi_0\rangle
\]
\[
E[\psi] = \frac{\langle \psi|H|\psi\rangle}{\langle \psi|\psi\rangle}
\]
\begin{itemize}
\item Central problem in quantum many-body physics
\end{itemize}
\end{frame}

%------------------------------------------------

\begin{frame}{Reference State}
\begin{itemize}
\item Start from Hartree–Fock determinant:
\[
|\phi_0\rangle
\]
\item Build correlations via excitation operators
\end{itemize}
\end{frame}

%================================================
\section{Second Quantization and Excitations}
%================================================

\begin{frame}{Fermionic Operators}
\[
\{a_p, a_q^\dagger\} = \delta_{pq}
\]
\begin{itemize}
\item Creation and annihilation operators
\end{itemize}
\end{frame}

%------------------------------------------------

\begin{frame}{Excitation Operators}
\[
T_1 = \sum_{ai} t_i^a a_a^\dagger a_i
\]
\[
T_2 = \sum_{abij} t_{ij}^{ab} a_a^\dagger a_b^\dagger a_j a_i
\]
\end{frame}

%================================================
\section{Jordan–Wigner Mapping}
%================================================

\begin{frame}{Fermion-to-Qubit Mapping}
\[
a_p = \frac{1}{2}(X_p + iY_p)\prod_{j<p} Z_j
\]
\[
a_p^\dagger = \frac{1}{2}(X_p - iY_p)\prod_{j<p} Z_j
\]
\end{frame}

%------------------------------------------------

\begin{frame}{Example Mapping}
\[
a_p^\dagger a_q \rightarrow \text{Pauli strings}
\]
\begin{itemize}
\item Leads to nonlocal strings of $Z$ operators
\end{itemize}
\end{frame}

%================================================
\section{Unitary Coupled Cluster}
%================================================

\begin{frame}{UCC Ansatz}
\[
|\psi\rangle = e^{T - T^\dagger} |\phi_0\rangle
\]
\end{frame}

%------------------------------------------------

\begin{frame}{BCH Expansion}
\[
e^{-T}He^{T} = H + [H,T] + \frac{1}{2}[[H,T],T] + \cdots
\]
\end{frame}

%------------------------------------------------

\begin{frame}{Classical CC vs UCC}
\begin{itemize}
\item CC:
\[
e^{T}
\]
(non-unitary)
\item UCC:
\[
e^{T - T^\dagger}
\]
(unitary)
\end{itemize}
\end{frame}

%================================================
\section{ADAPT-VQE Framework}
%================================================

\begin{frame}{Adaptive Ansatz}
\[
U(\theta) = \prod_k e^{\theta_k A_k}
\]
\end{frame}

%------------------------------------------------

\begin{frame}{Operator Pool}
\begin{itemize}
\item Fermionic excitations mapped to qubits
\item Anti-Hermitian operators:
\[
A_k = \tau_k - \tau_k^\dagger
\]
\end{itemize}
\end{frame}

%================================================
\section{Gradient Derivation via BCH}
%================================================

\begin{frame}{Energy Functional}
\[
E(\theta) = \langle \psi(\theta)|H|\psi(\theta)\rangle
\]
\end{frame}

%------------------------------------------------

\begin{frame}{Single Parameter Variation}
\[
|\psi(\theta)\rangle = e^{\theta A}|\psi\rangle
\]
\end{frame}

%------------------------------------------------

\begin{frame}{Derivative of Energy}
\[
\frac{dE}{d\theta}
=
\langle \psi | A^\dagger H + H A | \psi \rangle
\]
\end{frame}

%------------------------------------------------

\begin{frame}{Using Anti-Hermiticity}
\[
A^\dagger = -A
\]
\[
\frac{dE}{d\theta}
=
\langle \psi | [H,A] | \psi \rangle
\]
\end{frame}

%------------------------------------------------

\begin{frame}{BCH Interpretation}
\[
e^{-\theta A} H e^{\theta A}
=
H + \theta[H,A] + \frac{\theta^2}{2}[[H,A],A] + \cdots
\]
\begin{itemize}
\item Gradient corresponds to first-order term
\end{itemize}
\end{frame}

%================================================
\section{ADAPT Selection Rule}
%================================================

\begin{frame}{Gradient Norm}
\[
g_k = |\langle \psi | [H,A_k] | \psi \rangle|
\]
\end{frame}

%------------------------------------------------

\begin{frame}{Physical Meaning}
\begin{itemize}
\item Measures residual correlation
\item Analogous to CC amplitude equations
\end{itemize}
\end{frame}

%================================================
\section{Comparison with Coupled Cluster}
%================================================

\begin{frame}{CC Amplitude Equations}
\[
\langle \phi_\mu | e^{-T}He^T | \phi_0 \rangle = 0
\]
\end{frame}

%------------------------------------------------

\begin{frame}{ADAPT vs CC}
\begin{itemize}
\item CC:
\begin{itemize}
\item Solve nonlinear equations
\end{itemize}
\item ADAPT:
\begin{itemize}
\item Greedy operator selection
\end{itemize}
\end{itemize}
\end{frame}

%------------------------------------------------

\begin{frame}{ADAPT vs UCCSD}
\begin{itemize}
\item UCCSD:
\begin{itemize}
\item Fixed ansatz
\end{itemize}
\item ADAPT:
\begin{itemize}
\item Dynamically constructed
\item Typically fewer parameters
\end{itemize}
\end{itemize}
\end{frame}

%================================================
\section{Circuit Implementation}
%================================================

\begin{frame}{Operator Circuit}
\[
\begin{quantikz}
\lstick{} & \gate{e^{\theta_k A_k}} & \qw
\end{quantikz}
\]
\end{frame}

%------------------------------------------------

\begin{frame}{Ansatz Growth}
\[
U = e^{\theta_n A_n}\cdots e^{\theta_1 A_1}
\]
\end{frame}

%================================================
\section{Measurement Considerations}
%================================================

\begin{frame}{Commutator Evaluation}
\[
[H,A_k] = \sum_i c_i P_i
\]
\begin{itemize}
\item Measured via Pauli decomposition
\end{itemize}
\end{frame}

%------------------------------------------------

\begin{frame}{Cost Scaling}
\begin{itemize}
\item Scales with operator pool
\item Main bottleneck
\end{itemize}
\end{frame}

%================================================
\section{Convergence and Scaling}
%================================================

\begin{frame}{Stopping Criterion}
\[
\max_k g_k < \epsilon
\]
\end{frame}

%------------------------------------------------

\begin{frame}{Efficiency}
\begin{itemize}
\item Compact ansatz
\item Often faster convergence than UCCSD
\end{itemize}
\end{frame}

%================================================
\section{Interpretation}
%================================================

\begin{frame}{Many-Body View}
\begin{itemize}
\item ADAPT = iterative correlation builder
\item Similar to perturbation expansion
\end{itemize}
\end{frame}

%------------------------------------------------

\begin{frame}{Connection to BCH}
\begin{itemize}
\item Expansion defines effective Hamiltonian
\item Gradient = leading correction
\end{itemize}
\end{frame}

%================================================
\section{Summary}
%================================================

\begin{frame}{Summary}
\begin{itemize}
\item ADAPT-VQE builds ansatz dynamically
\item Gradient derived from commutator:
\[
\langle [H,A] \rangle
\]
\item Strong connection to:
\begin{itemize}
\item Coupled cluster theory
\item BCH expansion
\item Many-body perturbation theory
\end{itemize}
\end{itemize}
\end{frame}


%================================================
\section{Higher-Order BCH and Many-Body Structure}
%================================================

\begin{frame}{Baker--Campbell--Hausdorff Expansion}
For an operator $A$,
\[
\bar{H} = e^{-A} H e^{A}
\]
expands as
\[
\bar{H} = H + [H,A] + \frac{1}{2}[[H,A],A]
+ \frac{1}{6}[[[H,A],A],A] + \cdots
\]
\end{frame}

%------------------------------------------------

\begin{frame}{Nested Commutators Structure}
Define:
\[
\text{ad}_A(H) = [H,A]
\]
Then:
\[
\bar{H} = \sum_{n=0}^\infty \frac{1}{n!} \text{ad}_A^n(H)
\]

\begin{itemize}
\item Generates hierarchy of many-body terms
\item Higher commutators → higher excitation rank
\end{itemize}
\end{frame}

%------------------------------------------------

\begin{frame}{Example: Two-Body Hamiltonian}
Let
\[
H = H_1 + H_2
\]
\begin{itemize}
\item $H_1$: one-body
\item $H_2$: two-body
\end{itemize}

Then:
\[
[H,T_1] \rightarrow \text{one- and two-body terms}
\]
\[
[[H,T_1],T_1] \rightarrow \text{two- and three-body terms}
\]
\end{frame}

%------------------------------------------------

\begin{frame}{Physical Meaning of BCH}
\begin{itemize}
\item BCH builds effective Hamiltonian
\item Encodes correlations iteratively
\item ADAPT selects dominant directions in this expansion
\end{itemize}
\end{frame}

%================================================
\section{Diagrammatic Interpretation}
%================================================

\begin{frame}{Goldstone Diagram View}
\begin{itemize}
\item Each commutator corresponds to diagrams
\item $[H,T]$:
\begin{itemize}
\item Single interaction insertion
\end{itemize}
\item $[[H,T],T]$:
\begin{itemize}
\item Two interaction insertions
\end{itemize}
\end{itemize}
\end{frame}

%------------------------------------------------

\begin{frame}{Linked Diagram Theorem}
\begin{itemize}
\item Only connected diagrams contribute
\item Ensures size extensivity
\end{itemize}

\[
E = \langle \phi_0 | \bar{H} | \phi_0 \rangle
\]
\end{frame}

%------------------------------------------------

\begin{frame}{ADAPT-VQE Interpretation}
\begin{itemize}
\item Each selected operator adds diagram class
\item Greedy selection = dominant diagram contributions
\item Equivalent to truncated linked expansion
\end{itemize}
\end{frame}

%================================================
\section{Connection to MBPT}
%================================================

\begin{frame}{Many-Body Perturbation Theory}
\begin{itemize}
\item Expand in interaction strength:
\[
E = E^{(0)} + E^{(1)} + E^{(2)} + \cdots
\]
\end{itemize}
\end{frame}

%------------------------------------------------

\begin{frame}{MBPT vs Coupled Cluster}
\begin{itemize}
\item MBPT:
\begin{itemize}
\item Order-by-order expansion
\end{itemize}
\item CC:
\begin{itemize}
\item Infinite resummation
\end{itemize}
\end{itemize}
\end{frame}

%------------------------------------------------

\begin{frame}{ADAPT as Adaptive MBPT}
\begin{itemize}
\item ADAPT selects most important excitation
\item Equivalent to adaptive resummation
\item Avoids fixed truncation (like CCSD)
\end{itemize}
\end{frame}

%================================================
\section{Time-Dependent Coupled Cluster}
%================================================

\begin{frame}{Time-Dependent Ansatz}
\[
|\psi(t)\rangle = e^{T(t)} |\phi_0\rangle
\]
\end{frame}

%------------------------------------------------

\begin{frame}{TDCC Equation}
\[
i \frac{d}{dt}|\psi\rangle = H |\psi\rangle
\]
leads to
\[
i \dot{T} = \text{projection of } e^{-T}He^T
\]
\end{frame}

%------------------------------------------------

\begin{frame}{Linearized TDCC}
Small fluctuations:
\[
T(t) = T_0 + \delta T(t)
\]
leads to
\[
i \frac{d}{dt} \delta T = \mathcal{L} \delta T
\]
\end{frame}

%================================================
\section{Linear Response and EOM-CC}
%================================================

\begin{frame}{Equation-of-Motion CC}
\[
H e^T |R\rangle = \omega e^T |R\rangle
\]
\end{frame}

%------------------------------------------------

\begin{frame}{Linear Response Equation}
\[
(\omega I - \mathcal{L}) x = b
\]
\begin{itemize}
\item Same structure as HHL and TDHF
\end{itemize}
\end{frame}

%------------------------------------------------

\begin{frame}{RPA Structure}
\[
\begin{pmatrix}
A & B \\
-B^* & -A^*
\end{pmatrix}
\begin{pmatrix}
X \\ Y
\end{pmatrix}
=
\omega
\begin{pmatrix}
X \\ Y
\end{pmatrix}
\]
\end{frame}

%================================================
\section{Unified Picture}
%================================================

\begin{frame}{Unifying View}
\begin{itemize}
\item BCH → effective Hamiltonian
\item CC → resummation of diagrams
\item ADAPT → adaptive selection
\item TDCC → time evolution
\item Linear response → inverse operator
\end{itemize}
\end{frame}

%------------------------------------------------

\begin{frame}{Key Insight}
\begin{block}{Core idea}
ADAPT-VQE constructs a wavefunction by selecting operators that maximize:
\[
\langle [H,A_k] \rangle
\]
\end{block}

\begin{block}{Interpretation}
This is equivalent to selecting directions that contribute most strongly in:
\begin{itemize}
\item BCH expansion
\item diagrammatic expansion
\item perturbative corrections
\end{itemize}
\end{block}
\end{frame}

%------------------------------------------------

\begin{frame}{Final Conceptual Bridge}
\begin{itemize}
\item ADAPT-VQE = variational CC
\item Gradient = BCH first-order term
\item Ansatz = dynamically selected excitation algebra
\item Measurement = projection of many-body residual
\end{itemize}
\end{frame}


\end{document}
